\chapter{Implementation}

\section{Interpreter Core}
The FTCL implementation is mainly header-based and located under the \texttt{include} directory. The interpreter is built as a CMake interface library, while tests and benchmark executables are placed under the \texttt{test} directory. Optional CUDA support is enabled by the \texttt{FTCL\_ENABLE\_CUDA} CMake option.

Value representation follows the idea that script values are strings but may also cache structured internal data. For example, a value can be constructed from an integer, floating-point number, Boolean, list, dictionary, or variable name. When the string representation is requested, it is generated lazily. This reduces repeated formatting while preserving a string-oriented script interface.

\section{Expression Parser}
The expression parser implements a recursive-descent precedence hierarchy. The top-level function parses conditional expressions. It then descends through logical OR, logical AND, bitwise operators, equality, relational operators, shifts, additive operators, multiplicative operators, unary operators, and primary expressions.

For example, \texttt{a || b \&\& c} is parsed by the logical OR layer. The left operand is parsed as a logical AND expression. After the OR operator is matched, the right operand is again parsed as a logical AND expression, so \texttt{b \&\& c} is grouped before the OR operation is applied.

The parser also supports short-circuit evaluation. If the left operand of \texttt{||} is true, the right operand is parsed but not evaluated. This prevents side effects in skipped branches while preserving syntax validation and parser position.

\begin{algorithm}[H]
\caption{Expression Evaluation and Conditional Operator}
\label{alg:expr-eval}
\KwInput{interpreter $interp$, expression string $input$}
\KwOutput{expression value}
\Proc{\textsc{EvalExpr}($interp$, $input$)}{
  $parser \leftarrow \textsc{ExprParser}(interp, input)$\;
  $value \leftarrow parser.parseConditional()$\;
  $parser.skipWhitespace()$\;
  \If{not $parser.atEnd()$}{
    raise ``unexpected token in expression''\;
  }
  \KwReturn{$value$}\;
}
\Proc{\textsc{ParseConditional}($parser$)}{
  $condition \leftarrow \textsc{ParseLogicalOr}(parser)$\;
  \If{not $parser.match(\texttt{?})$}{
    \KwReturn{$condition$}\;
  }
  $takeTrue \leftarrow condition.isTrue()$\;
  $parser.setEvaluationEnabled(takeTrue)$\;
  $trueValue \leftarrow \textsc{ParseLogicalOr}(parser)$\;
  $parser.expect(\texttt{:})$\;
  $parser.setEvaluationEnabled(\neg takeTrue)$\;
  $falseValue \leftarrow \textsc{ParseConditional}(parser)$\;
  $parser.restoreEvaluationMode()$\;
  \uIf{$takeTrue$}{
    \KwReturn{$trueValue$}\;
  }
  \Else{
    \KwReturn{$falseValue$}\;
  }
}
\end{algorithm}

\begin{algorithm}[H]
\caption{Logical OR Short-circuit Evaluation}
\label{alg:expr-or}
\KwInput{expression parser $parser$}
\KwOutput{logical OR result}
\Proc{\textsc{ParseLogicalOr}($parser$)}{
  $lhs \leftarrow \textsc{ParseLogicalAnd}(parser)$\;
  \While{$parser.match(\texttt{||})$}{
    \uIf{$lhs.isTrue()$}{
      $parser.setEvaluationEnabled(false)$\;
      \textsc{ParseLogicalAnd}($parser$)\tcp*[r]{parse only; no side effects}
      $parser.restoreEvaluationMode()$\;
      $lhs \leftarrow 1$\;
    }
    \Else{
      $rhs \leftarrow \textsc{ParseLogicalAnd}(parser)$\;
      $lhs \leftarrow bool(lhs) \lor bool(rhs)$\;
    }
  }
  \KwReturn{$lhs$}\;
}
\end{algorithm}

\section{UVec Implementation}
UVec requires its element type to be trivially copyable. This is necessary because CPU/CUDA transfers are implemented as byte copies. The main supporting types are \texttt{Device}, \texttt{RawUPtr<T>}, \texttt{RawUMutPtr<T>}, \texttt{UniversalCopy}, and \texttt{Zeroable}.

The script-level \texttt{uvec} command manages UVec handles. It can create vectors from lists, create filled or zeroed vectors, get or set elements, return vector length, synchronize to a device, and convert a vector back to a script list. The command manager stores handle-to-UVec mappings under a mutex so that scripts can pass handles safely between commands.

\begin{algorithm}[H]
\caption{UVec Read Synchronization}
\label{alg:uvec-read}
\KwInput{UVec object $uvec$, target device $device$}
\KwOutput{read-only unified pointer}
\Proc{\textsc{AsUPtr}($uvec$, $device$)}{
  \uIf{$device = \mathrm{CPU}$}{
    \If{CPU buffer is not valid}{
      $source \leftarrow$ any valid CUDA buffer\;
      copy $source$ CUDA buffer to CPU buffer\;
      mark CPU buffer valid\;
    }
    \KwReturn{\textsc{RawUPtr}(CPU pointer, length, CPU)}\;
  }
  \Else{
    ensure CUDA buffer for $device$ exists\;
    \If{target CUDA buffer is not valid}{
      \uIf{CPU buffer is valid}{
        copy CPU buffer to target CUDA buffer\;
      }
      \Else{
        $source \leftarrow$ another valid CUDA buffer\;
        copy $source$ CUDA buffer to target CUDA buffer\;
      }
      mark target CUDA buffer valid\;
    }
    \KwReturn{\textsc{RawUPtr}(CUDA pointer, length, $device$)}\;
  }
}
\end{algorithm}

\begin{algorithm}[H]
\caption{UVec Mutable Pointer Access}
\label{alg:uvec-write}
\KwInput{UVec object $uvec$, target device $device$}
\KwOutput{mutable unified pointer}
\Proc{\textsc{AsMutUPtr}($uvec$, $device$)}{
  ensure storage exists on $device$\;
  \If{target device is not valid}{
    synchronize latest data to target device\;
  }
  \ForEach{$other \in$ devices and $other \ne device$}{
    mark $other$ invalid\;
  }
  mark $device$ valid and newest\;
  \KwReturn{\textsc{RawUMutPtr}(target pointer, length, $device$)}\;
}
\end{algorithm}

\section{CPU Geometry Algorithms}
The CPU geometry library is implemented in \texttt{include/geometry.hpp}. The basic types are Point, Segment, and BoundingBox. The library includes distance, squared distance, dot product, cross product, orientation, point-on-segment checks, segment intersection, line intersection, point-line distance, point-segment distance, polygon area, polygon perimeter, point-in-polygon classification, convex hull, bounding box, and closest-pair distance.

Batch algorithms use flattened arrays. A point set is stored as \(x_0,y_0,x_1,y_1,\ldots\). A segment set is stored as \(x_1,y_1,x_2,y_2\) per segment. An AABB set is stored as \(min_x,min_y,max_x,max_y\) per box. This format is easy for scripts, UVec, and CUDA kernels to share.

\begin{table}[H]
\centering
\caption{Representative geometry algorithms and asymptotic costs}
\label{tab:algorithm-complexity}
\footnotesize
\setlength{\tabcolsep}{3.5pt}
\begin{tabularx}{\textwidth}{@{}>{\ttfamily\raggedright\arraybackslash}p{0.31\textwidth} >{\raggedright\arraybackslash}X >{\centering\arraybackslash}p{0.12\textwidth} >{\raggedright\arraybackslash}X@{}}
\toprule
Command or function & Input scale & CPU complexity & GPU suitability \\
\midrule
distance & one point pair & \(O(1)\) & Low; scalar command \\
segment\_intersect & one segment pair & \(O(1)\) & Low; scalar command \\
polygon\_area & \(P\) polygon vertices & \(O(P)\) & Medium for batched polygons \\
convex\_hull & \(N\) points & \(O(N\log N)\) & Low in this project \\
bbox\_reduce & \(N\) points & \(O(N)\) & High with parallel reduction \\
centroid & \(N\) points & \(O(N)\) & High with parallel reduction \\
batch\_distance\_matrix & \(N\) by \(Q\) points & \(O(NQ)\) & Very high; independent outputs \\
nearest\_point & \(N\) data, \(Q\) queries & \(O(NQ)\) & High; independent queries \\
range\_count\_circle & \(N\) points, \(Q\) circles & \(O(NQ)\) & High; independent queries \\
batch\_segment\_intersect & \(K\) segment pairs & \(O(K)\) & High; independent pairs \\
collision\_aabb & \(K\) box pairs & \(O(K)\) & High; independent pairs \\
transform\_points & \(N\) points & \(O(N)\) & Very high; independent points \\
\bottomrule
\end{tabularx}
\end{table}

The table explains why not every geometry command is accelerated in the same way. Scalar commands are useful for script expressiveness but are too small to amortize CUDA launch overhead. Batch commands with \(O(NQ)\) or \(O(N)\) independent work items are more suitable for GPU execution. Algorithms such as convex hull require sorting and sequential stack-like structure, so they are kept as CPU-first algorithms in this thesis.

\section{CUDA Geometry Backend}
The CUDA backend is implemented in \texttt{src/geometry\_cuda.cu}.
It is declared in \texttt{include/geometry\_cuda.hpp}.
It supports data-parallel algorithms such as distance matrix, nearest point, circle range count, point transformation, batch orientation, AABB collision, polygon batch area, and distance to polyline.

For a distance matrix with \(N\) source points and \(Q\) query points, the output has \(NQ\) entries. A CUDA thread computes one output entry by loading one source point and one query point, computing the Euclidean distance, and writing the result. This mapping is simple and effective because each output element is independent.

\section{Script Interface}
The following example shows how a script can create two point sets, run a CUDA nearest-point query, and read back the result.

\begin{lstlisting}[language=bash,caption={FTCL geometry scripting example}]
set dataset [geom uvec_points {{0 0} {1 2} {3 4} {5 8}} cuda:0]
set query   [geom uvec_points {{2 2} {4 6}} cuda:0]
set result  [geom nearest_point $dataset $query cuda:0]
puts [uvec to_list $result]
\end{lstlisting}

The user does not manually allocate device memory. The script describes the task, UVec manages synchronization, and the geom command selects the backend.

\begin{algorithm}[H]
\caption{Geometry Command Dispatch for Nearest-point Query}
\label{alg:geom-dispatch}
\KwInput{interpreter $interp$, command arguments $argv$}
\KwOutput{UVec handle containing nearest-point results}
\Proc{\textsc{GeomNearestPoint}($interp$, $argv$)}{
  $datasetId \leftarrow \textsc{ParseHandle}(argv[2])$\;
  $queryId \leftarrow \textsc{ParseHandle}(argv[3])$\;
  $device \leftarrow \textsc{ParseOptionalDevice}(argv[4], \mathrm{CPU})$\;
  $dataset \leftarrow uvecManager.lookup(datasetId)$\;
  $queries \leftarrow uvecManager.lookup(queryId)$\;
  \uIf{$device$ is CUDA}{
    $datasetPtr \leftarrow dataset.asUPtr(device)$\;
    $queryPtr \leftarrow queries.asUPtr(device)$\;
    $output \leftarrow \textsc{UVec.Uninitialized}(queryCount \times 2, device)$\;
    $outputPtr \leftarrow output.asMutUPtr(device)$\;
    launch \textsc{NearestPointCuda}($outputPtr$, $datasetPtr$, $queryPtr$)\;
    \KwReturn{$uvecManager.create(output)$}\;
  }
  \Else{
    $datasetCpu \leftarrow dataset.cpuVector()$\;
    $queriesCpu \leftarrow queries.cpuVector()$\;
    $values \leftarrow \textsc{NearestPointCpu}(datasetCpu, queriesCpu)$\;
    $output \leftarrow \textsc{UVec.FromCpu}(values)$\;
    \KwReturn{$uvecManager.create(output)$}\;
  }
}
\end{algorithm}
